% Options for packages loaded elsewhere
\PassOptionsToPackage{unicode}{hyperref}
\PassOptionsToPackage{hyphens}{url}
\PassOptionsToPackage{dvipsnames,svgnames,x11names}{xcolor}
\documentclass[
  10pt,
  fontset=fandol]{ctexart}
\usepackage{xcolor}
\usepackage[margin=16mm]{geometry}
\usepackage{amsmath,amssymb}
\setcounter{secnumdepth}{-\maxdimen} % remove section numbering
\usepackage{iftex}
\ifPDFTeX
  \usepackage[T1]{fontenc}
  \usepackage[utf8]{inputenc}
  \usepackage{textcomp} % provide euro and other symbols
\else % if luatex or xetex
  \usepackage{unicode-math} % this also loads fontspec
  \defaultfontfeatures{Scale=MatchLowercase}
  \defaultfontfeatures[\rmfamily]{Ligatures=TeX,Scale=1}
\fi
\usepackage{lmodern}
\ifPDFTeX\else
  % xetex/luatex font selection
\fi
% Use upquote if available, for straight quotes in verbatim environments
\IfFileExists{upquote.sty}{\usepackage{upquote}}{}
\IfFileExists{microtype.sty}{% use microtype if available
  \usepackage[]{microtype}
  \UseMicrotypeSet[protrusion]{basicmath} % disable protrusion for tt fonts
}{}
\makeatletter
\@ifundefined{KOMAClassName}{% if non-KOMA class
  \IfFileExists{parskip.sty}{%
    \usepackage{parskip}
  }{% else
    \setlength{\parindent}{0pt}
    \setlength{\parskip}{6pt plus 2pt minus 1pt}}
}{% if KOMA class
  \KOMAoptions{parskip=half}}
\makeatother
\usepackage{longtable,booktabs,array}
\newcounter{none} % for unnumbered tables
\usepackage{calc} % for calculating minipage widths
% Correct order of tables after \paragraph or \subparagraph
\usepackage{etoolbox}
\makeatletter
\patchcmd\longtable{\par}{\if@noskipsec\mbox{}\fi\par}{}{}
\makeatother
% Allow footnotes in longtable head/foot
\IfFileExists{footnotehyper.sty}{\usepackage{footnotehyper}}{\usepackage{footnote}}
\makesavenoteenv{longtable}
\setlength{\emergencystretch}{3em} % prevent overfull lines
\providecommand{\tightlist}{%
  \setlength{\itemsep}{0pt}\setlength{\parskip}{0pt}}
\usepackage{amsmath,amssymb,mathtools,needspace}
\xeCJKDeclareCharClass{CJK}{"2032,"0400->"04FF,"2160->"216B,"2460->"2473,"25A0,"25C7,"25CB,"25CF}
\IfFontExistsTF{Arial Unicode MS}{\xeCJKsetup{AutoFallBack=true}\setCJKfallbackfamilyfont{\CJKrmdefault}{Arial Unicode MS}}{}
\setcounter{secnumdepth}{0}
\setlength{\emergencystretch}{2em}
\allowdisplaybreaks
\usepackage{bookmark}
\IfFileExists{xurl.sty}{\usepackage{xurl}}{} % add URL line breaks if available
\urlstyle{same}
\hypersetup{
  pdftitle={等距节点插值公式},
  colorlinks=true,
  linkcolor={Maroon},
  filecolor={Maroon},
  citecolor={Blue},
  urlcolor={Blue},
  pdfcreator={LaTeX via pandoc}}

\title{等距节点插值公式}
\author{}
\date{}

\begin{document}
\maketitle

\subsubsection{2.5.2 等距节点插值公式}\label{ux7b49ux8dddux8282ux70b9ux63d2ux503cux516cux5f0f}

将 Newton 差商插值多项式（2.4.6）中各阶差商用相应差分代替，就可得到各种形式的等距节点插值公式。这里只推导常用的前插公式与后插公式。

如果节点 \(x_k=x_0+kh\;(k=0,1,\cdots,n)\)，要计算 \(x_0\) 附近点 \(x\) 的函数 \(f(x)\) 的值，可令 \(x=x_0+th\)，\(0\le t\le1\)，于是

\[
\omega_{k+1}(x)=\prod_{j=0}^{k}(x-x_j)=t(t-1)\cdots(t-k)h^{k+1}.
\]

将此式及式（2.5.7）代入式（2.4.6），得

\[
N_n(x_0+th)=f_0+t\Delta f_0+\frac{t(t-1)}{2!}\Delta^2 f_0+\cdots+\frac{t(t-1)\cdots(t-n+1)}{n!}\Delta^n f_0.
\tag{2.5.10}
\]

式（2.5.10）称为 \textbf{Newton 前插公式}，其余项由式（2.2.14）得

\[
R_n(x)=\frac{t(t-1)\cdots(t-n)}{(n+1)!}h^{n+1}f^{(n+1)}(\xi),\qquad\xi\in(x_0,x_n).
\tag{2.5.11}
\]

如果要求表示函数在 \(x_n\) 附近的值 \(f(x)\)，此时应用 Newton 插值公式（2.4.6），插值点应按 \(x_n,x_{n-1},\cdots,x_0\) 的次序排列，有

\[
\begin{aligned}
N_n(x)={}&f(x_n)+f[x_n,x_{n-1}](x-x_n)+f[x_n,x_{n-1},x_{n-2}](x-x_n)(x-x_{n-1})\\
&+\cdots+f[x_n,x_{n-1},\cdots,x_0](x-x_n)\cdots(x-x_1).
\end{aligned}
\]

作变换 \(x=x_n+th\;(-1\le t\le0)\)，并利用式（2.5.8），代入上式得

\[
\begin{aligned}
N_n(x_n+th)={}&f_n+t\nabla f_n+\frac{t(t+1)}{2!}\nabla^2 f_n+\cdots\\
&+\frac{t(t+1)\cdots(t+n-1)}{n!}\nabla^n f_n.
\end{aligned}
\tag{2.5.12}
\]

式（2.5.12）称为 \textbf{Newton 后插公式}，其余项

\[
R_n(x)=f(x)-N_n(x_n+th)=\frac{t(t+1)\cdots(t+n)h^{n+1}f^{(n+1)}(\xi)}{(n+1)!},\qquad\xi\in(x_0,x_n).
\]

利用 Newton 前插公式（2.5.10）计算函数值时，其系数就是 \(f(x)\) 在 \(x_0\) 的各阶向前差分（见表 2.6）。若用 Newton 后插公式（2.5.12）求 \(f(x)\) 的值，因 \(x\) 在 \(x_n\) 附近，则其系数为 \(f(x)\) 在点 \(x_n\) 的各阶向后差分。

等距节点插值公式有不少实际应用，例如，很多工程设计计算都需要查各种函数表，用计算机计算时就必须解决计算机查表问题。如果把整个函数表存入内存，往往占用单元太多；如果用一个解析表达式近似该函数，又可能达不到精度要求。因此，采用存放大间隔函数表，并用插值公式计算函数近似值，是一种可行的方案。

\textbf{例 2.4} 在微电机设计计算中需要查磁化曲线表，通常给出的表（见表 2.7）是磁密 \(B\) 每间隔 100 高斯磁路每厘米长所需安匝数 \(at\) 的值，下面要解决 \(B\) 从 \(4\,000\) 至

\begin{quote}
例 2.4 的区间上限及解答续 PDF 42。
\end{quote}

\href{../../source-images/pdf-042.jpeg}{原页：书页 29／PDF 42}

\begin{quote}
页首接 PDF 41 例 2.4 的区间；页末 2.6 首句续 PDF 43。
\end{quote}

\(11\,000\) 区间的查表问题。

\textbf{解} 为节省计算机存储单元，采用每 500 高斯存入一个 \(at\) 值。为了分析使用几阶插值公式合适，应先列出差分表。表 2.7 是以硅钢片的磁化曲线为例所造的差分表。

从差分表中看到三阶差分近似于 0，因此计算时只需用二阶差分，当 \(4\,000\le B\le10\,500\) 时用 Newton 前插公式，当 \(10\,500<B\le11\,000\) 时用 Newton 后插公式。例如，求 \(f(5\,200)\) 时取 \(B_0=5\,000\)，\(f_0=1.58\)，\(\Delta f_0=0.11\)，\(\Delta^2 f_0=0.01\)，\(h=500\)，\(B=5\,200\)，\(t=0.4\)，于是由式（2.5.10），取 \(n=2\)，得

\[
f(5\,200)\approx1.58+0.4\times0.11+\frac{0.4\times(-0.6)}{2}\times0.01\approx1.62.
\]

这个结果与直接查表得到的值相同，说明用此法在计算机上求 \(at\) 值是可行的。具体的查表程序作为练习由读者完成。

\textbf{表 2.7}

{\def\LTcaptype{none} % do not increment counter
\begin{longtable}[]{@{}
  >{\raggedleft\arraybackslash}p{(\linewidth - 10\tabcolsep) * \real{0.1667}}
  >{\raggedleft\arraybackslash}p{(\linewidth - 10\tabcolsep) * \real{0.1667}}
  >{\raggedleft\arraybackslash}p{(\linewidth - 10\tabcolsep) * \real{0.1667}}
  >{\raggedleft\arraybackslash}p{(\linewidth - 10\tabcolsep) * \real{0.1667}}
  >{\raggedleft\arraybackslash}p{(\linewidth - 10\tabcolsep) * \real{0.1667}}
  >{\raggedleft\arraybackslash}p{(\linewidth - 10\tabcolsep) * \real{0.1667}}@{}}
\toprule\noalign{}
\begin{minipage}[b]{\linewidth}\raggedleft
\(k\)
\end{minipage} & \begin{minipage}[b]{\linewidth}\raggedleft
\(B_k\)
\end{minipage} & \begin{minipage}[b]{\linewidth}\raggedleft
\(at_k=f(B_k)\)
\end{minipage} & \begin{minipage}[b]{\linewidth}\raggedleft
\(\Delta f_k\)
\end{minipage} & \begin{minipage}[b]{\linewidth}\raggedleft
\(\Delta^2 f_k\)
\end{minipage} & \begin{minipage}[b]{\linewidth}\raggedleft
\(\Delta^3 f_k\)
\end{minipage} \\
\midrule\noalign{}
\endhead
\bottomrule\noalign{}
\endlastfoot
\(0\) & \(4\,000\) & \(1.38\) & \(0.10\) & \(0\) & \(0.01\) \\
\(1\) & \(4\,500\) & \(1.48\) & \(0.10\) & \(0.01\) & \(0\) \\
\(2\) & \(5\,000\) & \(1.58\) & \(0.11\) & \(0.01\) & \(0\) \\
\(3\) & \(5\,500\) & \(1.69\) & \(0.12\) & \(0.01\) & \(0.02\) \\
\(4\) & \(6\,000\) & \(1.81\) & \(0.13\) & \(0.03\) & \(-0.01\) \\
\(5\) & \(6\,500\) & \(1.94\) & \(0.16\) & \(0.02\) & \(0.02\) \\
\(6\) & \(7\,000\) & \(2.10\) & \(0.18\) & \(0.04\) & \(0\) \\
\(7\) & \(7\,500\) & \(2.28\) & \(0.22\) & \(0.04\) & \(0\) \\
\(8\) & \(8\,000\) & \(2.50\) & \(0.26\) & \(0.04\) & \(0.01\) \\
\(9\) & \(8\,500\) & \(2.76\) & \(0.30\) & \(0.05\) & \(0.02\) \\
\(10\) & \(9\,000\) & \(3.06\) & \(0.35\) & \(0.07\) & \(0.01\) \\
\(11\) & \(9\,500\) & \(3.41\) & \(0.42\) & \(0.08\) & \(0.02\) \\
\(12\) & \(10\,000\) & \(3.83\) & \(0.50\) & \(0.10\) & \\
\(13\) & \(10\,500\) & \(4.33\) & \(0.60\) & & \\
\(14\) & \(11\,000\) & \(4.93\) & & & \\
\end{longtable}
}

\end{document}
